\emph{(i).} Under Hypothesis (S) (\claimref{standard-model.hyp-s-slot-statistics}) the generators of $\su(3)$ and $\su(2)_{L}$ (\claimref{standard-model.automorphisms-preserve-slot-pair-split-are}) preserve $j$ and $k$, and each species, the states of given $(j, k)$, is irreducible under $\mathfrak{g}$: $\su(3)$ acts irreducibly on $\Lambda^{k}\C^{3}$ and $\su(2)_{L}$ on $\Lambda^{j}\C^{2}$. The six species have dimension and hypercharge $(6, \tfrac16)$, $(3, -\tfrac23)$, $(3, \tfrac13)$, $(2, -\tfrac12)$, $(1, 1)$, $(1, 0)$, pairwise distinct, so they are the isotypic components of the sixteen. If $g\mathfrak{g}g^{-1} = \mathfrak{g}$ then $g$ carries each isotypic component to an isotypic component of the same dimension, and $gYg^{-1}$ lies in the centre of $\mathfrak{g}$, so $gYg^{-1} = cY$. The species $Q$ is the only one of dimension $6$, so $g$ preserves it, and comparing the value of $Y$ on it gives $c = 1$. So $g$ commutes with $Y$, and since the six values of $Y$ are distinct $g$ preserves every species. An operator commuting with $\mathfrak{g}$ is a scalar on each species, so $g$ commutes with it; $n_{B}$ is such an operator (\claimref{standard-model.proton-decay-only-through-slot-pair}). With a family space the isotypic components are the species tensored with it, of dimensions in the same ratio, and the argument is unchanged.

\emph{(ii).} $X_{A}$ is the automorphism of the algebra of the ten induced by the reflection of the generators $v_{a} - w_{a}$, $a \in A$; on the module it is implemented, up to the grading, by the product of the $v_{a} + w_{a}$, $a \in A$, which sends $e_{S}$ to $\pm e_{S \triangle A}$. An operator that is a function $\phi$ of the occupation is carried to $\phi(S \triangle A)$. The check runs over the $32$ subsets: $n_{B}(S \triangle A) = -n_{B}(S)$ for all $S$ exactly when $A \supseteq \{1, 2, 3\}$, where $k \mapsto 3 - k$. For the normaliser: if $a \in A$ and $b \notin A$ are both pair modes, $X_{A}(v_{a}w_{b}) = w_{a}w_{b}$ changes $k$ and is not in $\mathfrak{g}$, and likewise for the slot modes with $v_{4}w_{5}$; so $A$ contains all or none of the pair modes and all or none of the slot modes. Then $X_{A}$ sends $T_{3R} = \tfrac12(N_{4} + N_{5} - 1)$ to $\pm T_{3R}$ and $B - L$ to $\pm(B - L)$ with independent signs, and $Y = T_{3R} + \tfrac12(B - L)$ to a multiple of $Y$ only when the signs agree: $A = \emptyset$ or $A = \{1, \dots, 5\}$. The total exchange does normalise $\mathfrak{g}$: $v_{a}w_{b} \mapsto w_{a}v_{b} = -v_{b}w_{a}$ for $a \neq b$, and $Y \mapsto -Y$. It changes the parity of $|S|$, so it maps the sixteen to the odd states. The odd states carry the weights opposite to those of the sixteen, and $\mathfrak{g}$ is a compact form acting by anti-Hermitian generators, so the odd states are the complex conjugate of the sixteen as a module; composing gives an antilinear map of the sixteen to itself normalising $\mathfrak{g}$. For an antilinear $g$ normalising $\mathfrak{g}$ the argument of (i) applies to $g$ with $gYg^{-1} = cY$, $c$ real, and gives $c = 1$ and $g\,n_{B}\,g^{-1} = n_{B}$; $g$ anticommutes with the unit, so it conjugates the rotation generated by the unit times $n_{B}$ to its inverse.

\emph{(iii).} $K = [H_{u}, H_{d}]$ satisfies $K^{\dagger} = -K$. For a $3 \times 3$ matrix $\det(-K) = -\det K$, so $(\det K)^{*} = \det K^{*} = \det (K^{\dagger})^{T} = \det K^{\dagger} = \det(-K) = -\det K$. $[UH_{u}U^{\dagger}, UH_{d}U^{\dagger}] = UKU^{\dagger}$ has the determinant of $K$. $[H_{u}^{*}, H_{d}^{*}] = K^{*}$, with determinant $(\det K)^{*} = -\det K$. If both are real symmetric, $K$ is real and antisymmetric, and $\det K = \det K^{T} = \det(-K) = -\det K$, so $\det K = 0$. In Fact \claimref{synthesis.imported-jarlskog-commutator-identity}, $\det[H_{u}, H_{d}] = 2I_{4}J F_{u}F_{d}$ with $F_{u}$, $F_{d}$ functions of the eigenvalues, which are real and unchanged by conjugation and by a change of basis, and nonzero for nondegenerate masses; so $J$ changes as $D$ does divided by $I_{4}$: it is fixed by a change of basis and reversed by conjugation.

\emph{(iv).} A linear map of the class fixes the baryon charge by (i) and $J$ by (iii); an antilinear one reverses the baryon charge by (ii) and $J$ by (iii). An automorphism of the laws in the sense stated fixes $J$, so with $J \neq 0$ it is linear and fixes the baryon charge. If $H_{u}$, $H_{d}$ are real in a family basis, conjugation in that basis fixes them and, composed with the total exchange, reverses the baryon charge. This clause is argued and not proved in one respect: that the maps relating a state to its charge conjugate at a fixed direction of time are exhausted by this class is the standard reading of charge conjugation with parity on left-handed fields, taken from the literature and not derived here.

\emph{Reading.} The properties of $C$ are those stated in \claimref{dark-matter.orientation-arrow-seventeenth-refutation}; the phase requirement is that of \claimref{synthesis.cabibbo-from-one-type-quark-stages-selects-near-quarter-turn-phase}; the identification of the frequency sign with the direction of time for a fixed law is \claimref{one-relation.arrow-frequency-are-one-bit-level}. The two readings are the two ways an antilinear anti-automorphism can act on a term and its adjoint; nothing stored selects one.
